
\section{定界余链与 $A_\infty$ Maurer--Cartan 方程}
\label{sec:3-6}

到目前为止，我们已经在不假设任何无障碍条件的情况下，为相对 spin
拉格朗日子流形构造了滤过 $A_\infty$ 代数。由于 $\mathfrak m_0$ 的存在，
算子 $\mathfrak m_1$ 一般不能定义上同调。本节研究把 $\mathfrak m_1$
形变成能够定义上同调的算子时所遇到的障碍。为此，我们先在第 3.6.1 小节
引入滤过 $A_\infty$ 代数中定界余链（Maurer--Cartan 方程的解）的代数概念；
再在第 3.6.2 小节构造一列经典上同调类，用来描述定界余链存在性的障碍。
进一步，第 3.6.3 小节将利用 $A_\infty$ 代数的幺元或同伦幺元，引入一个
较弱的无障碍条件。这些构造都可以完全用代数方式完成。

\subsection{定界余链与形变}
\label{subsec:3-6-1}

令 $(C,\mathfrak m)$ 是 $\Lambda_{0,\mathrm{nov}}$ 上的滤过
$A_\infty$ 代数，并取
$b\in(C[1])^0=C^1$，满足
$b\equiv0\pmod{\Lambda^+_{0,\mathrm{nov}}}$。置
\begin{equation}
  e^b:=1+b+b\otimes b+\cdots
  +\underbrace{b\otimes\cdots\otimes b}_{k\text{ 个}}+\cdots
  \in\widehat B(C[1]).
  \tag{3.6.1}\label{eq:3-6-1}
\end{equation}
这里与通常的指数函数定义不同，没有加入阶乘。这是因为第 3.5 节在几何上
实现映射 $\mathfrak m_{k,\beta}$ 时，我们在 $\mathcal M_{k+1}$ 的
$k!$ 个分支中只采用了主分支。条件
$b\equiv0\pmod{\Lambda^+_{0,\mathrm{nov}}}$ 蕴含
\eqref{eq:3-6-1} 右端关于能量滤过所定义的拓扑收敛，因而确实给出
$\widehat B(C[1])$ 的一个元素。

下面记录一个事实。

\begin{lemma}{3.6.2}\label{lem:3-6-2}
对非零元素 $x\in\widehat B(C[1])$，存在
$b\in B_1(C[1])=C[1]$ 使 $x=e^b$，当且仅当
\[
  \Delta x=x\mathbin{\overline\otimes}x,
\]
其中 $\Delta$ 是由 (3.2.14) 定义在 $\widehat B(C[1])$ 上的余积。
\end{lemma}

\begin{proof}
“仅当”方向容易直接验证，所以只证“若”方向。写
$x=\sum_{k=0}^{\infty}x_k$，其中
$x_k\in\widehat B_k(C[1])$，并假设
\begin{equation}
  \Delta x=x\mathbin{\overline\otimes}x.
  \tag{3.6.3}\label{eq:3-6-3}
\end{equation}
注意
\[
  \Delta\colon\widehat B(C[1])\longrightarrow
  \widehat B(C[1])\mathbin{\overline\otimes}\widehat B(C[1]);
\]
这里以 $\mathbin{\overline\otimes}$ 表示上式右端的张量积。
我们把
\[
  \bigoplus_{k_1+k_2=k}
  \widehat B_{k_1}(C[1])
  \mathbin{\overline\otimes}\widehat B_{k_2}(C[1])
\]
在
$\widehat B(C[1])\mathbin{\overline\otimes}\widehat B(C[1])$
中的分量称为 $k$ 阶项，并把 $y$ 的 $k$ 阶项记作 $(y)_k$。
同样，也把 $y\in\widehat B(C[1])$ 的
$\widehat B_k(C[1])$-分量记作 $(y)_k$。

利用余积 $\Delta$ 的定义 (3.2.14)，比较~\eqref{eq:3-6-3} 两边的零阶项，
得到
\[
  (x_0)^2=x_0,\qquad x_0\in\Lambda_{0,\mathrm{nov}},
\]
故 $x_0=0$ 或 $1$。若 $x_0=1$，比较~\eqref{eq:3-6-3} 的一阶项得到
\[
  1\mathbin{\overline\otimes}x_1
  +x_1\mathbin{\overline\otimes}1
  =1\mathbin{\overline\otimes}x_1
  +x_1\mathbin{\overline\otimes}1.
\]
换言之，这个方程对一阶项没有施加限制。置 $b:=x_1$。一般地，对二阶项，
容易看出余积满足
\[
  (\Delta x_2)_2
  =1\mathbin{\overline\otimes}x_2+\pi(x_2)
   +x_2\mathbin{\overline\otimes}1.
\]
这里
\[
  \pi\colon\widehat B_2(C[1])\longrightarrow
  \widehat B_1(C[1])\mathbin{\overline\otimes}\widehat B_1(C[1])
\]
由 $\pi(x\otimes y)=x\mathbin{\overline\otimes}y$ 定义。因此，
\eqref{eq:3-6-3} 的二阶项变成
\[
  1\mathbin{\overline\otimes}x_2+\pi(x_2)
  +x_2\mathbin{\overline\otimes}1
  =1\mathbin{\overline\otimes}x_2
  +x_1\mathbin{\overline\otimes}x_1
  +x_2\mathbin{\overline\otimes}1,
\]
所以 $x_2=x_1\otimes x_1=b\otimes b$。逐阶重复比较
\eqref{eq:3-6-3} 两边，便证明 $x=e^b$。若 $x_0=0$，同样的论证给出
$x=0$；这与假设相斥。引理得证。
\end{proof}

\begin{definition}{3.6.4}\label{def:3-6-4}
设 $b\in(C[1])^0$ 且
$b\equiv0\pmod{\Lambda^+_{0,\mathrm{nov}}}$。若
\[
  \widehat d(e^b)=0,
\]
就称 $b$ 是一个\emph{定界余链}\liftIndex{I004}\liftIndex{I009}，
或称它是 Maurer--Cartan 方程\liftIndex{I068}\liftIndex{I078}的一个解。
（回忆 $\widehat d=\sum\widehat{\mathfrak m}_k$，且
$\widehat{\mathfrak m}_k$ 由 (3.2.15) 定义。）若滤过 $A_\infty$ 代数
存在定界余链，就称它\emph{无障碍}\liftIndex{I069}\liftIndex{I110}；
否则称它\emph{有障碍}。
把所有定界余链 $b$ 构成的集合记作 $\widehat{\mathcal M}(C)$。
\end{definition}

\begin{remark}{3.6.5}\label{rem:3-6-5}
方程 $\widehat d(e^b)=0$ 是 Maurer--Cartan 方程或
Batalin--Vilkovisky 主方程的一个非齐次 $A_\infty$ 版本
\textbf{[BaVi93]}。事实上，它可以写成
\[
  \delta b_\beta+
  \sum_{\substack{k=0,1,2,\ldots\\
       \beta_{k+1},\beta_1,\ldots,\beta_k\in G\\
       \beta_{k+1}+\beta_1+\cdots+\beta_k=\beta\\
       (k,\beta_{k+1})\ne(1,(0,0))}}
  \mathfrak m_{k,\beta_{k+1}}
  (b_{\beta_1},\ldots,b_{\beta_k})=0,
\]
其中
\[
  b=\sum_\beta T^{\lambda(\beta)}e^{\mu(\beta)/2}b_\beta.
\]
BV 主方程与形变理论的关系见
\textbf{[ASKZ97, GugSt86, HiSc97, Kon03, Sche98]}；另见
\textbf{[Fuk02III]}。
\end{remark}

令 $(C,\mathfrak m)$ 是具有定界余链 $b_1,b_0$ 的滤过
$A_\infty$ 代数。如下定义映射
\[
  \delta_{b_1,b_0}\colon C[1]\longrightarrow C[1].
\]

\begin{definition}{3.6.6}\label{def:3-6-6}
\[
  \delta_{b_1,b_0}(x)
  =\sum_{k_1,k_0\geq0}
  \mathfrak m_{k_1+k_0+1}
  \bigl(
    \underbrace{b_1,\ldots,b_1}_{k_1},x,
    \underbrace{b_0,\ldots,b_0}_{k_0}
  \bigr).
\]
\end{definition}

在不会混淆的意义下，可以把这个定义简写成
$\delta_{b_1,b_0}(x)=\mathfrak m(e^{b_1}xe^{b_0})$。\mbox{于是}
\begin{equation}
\begin{aligned}
  \widehat d(e^{b_1}xe^{b_0})
  &={}e^{b_1}\delta_{b_1,b_0}(x)e^{b_0}
    +\widehat d(e^{b_1})xe^{b_0}\\
  &\quad{}+(-1)^{\deg x+1}e^{b_1}x\widehat d(e^{b_0}).
\end{aligned}
\tag{3.6.7}\label{eq:3-6-7}
\end{equation}
若 $\widehat d(e^{b_1})=\widehat d(e^{b_0})=0$，第二、三项便消失。

\begin{lemma}{3.6.8}\label{lem:3-6-8}
若 $\widehat d(e^{b_1})=\widehat d(e^{b_0})=0$，则
\[
  \delta_{b_1,b_0}\circ\delta_{b_1,b_0}=0.
\]
\end{lemma}

\begin{proof}
由~\eqref{eq:3-6-7}，
\[
\begin{aligned}
0
&=\widehat d\bigl(\widehat d(e^{b_1}xe^{b_0})\bigr)\\
&=\widehat d\bigl(e^{b_1}\delta_{b_1,b_0}(x)e^{b_0}\bigr)\\
&=e^{b_1}
  \bigl((\delta_{b_1,b_0}\circ\delta_{b_1,b_0})(x)\bigr)e^{b_0},
\end{aligned}
\]
结论立即成立。
\end{proof}

还可以对任意满足
$b\in C[1]^0$、$b\equiv0\pmod{\Lambda^+_{0,\mathrm{nov}}}$
的上链 $b$ 形变滤过 $A_\infty$ 代数的结构；这里并不要求 $b$ 是定界余链。
定义如下。以后将在第 5.2.2 小节讨论滤过 $A_\infty$ 同态的形变
\liftIndex{I080}。

\begin{definition}{3.6.9}\label{def:3-6-9}
对任意满足上述条件的上链 $b$，置
\[
\begin{aligned}
  \mathfrak m_k^b(x_1,\ldots,x_k)
  &={}
  \sum_{\ell_0,\ldots,\ell_k\geq0}
  \mathfrak m_{k+\sum_i\ell_i}
  \bigl(
    \underbrace{b,\ldots,b}_{\ell_0},x_1,
    \underbrace{b,\ldots,b}_{\ell_1},\ldots,\\[-0.2em]
  &\hspace{12em}
    \underbrace{b,\ldots,b}_{\ell_{k-1}},x_k,
    \underbrace{b,\ldots,b}_{\ell_k}
  \bigr)\\
  &={}
  \mathfrak m(e^b x_1e^b x_2\cdots x_{k-1}e^b x_ke^b),
\end{aligned}
\]
其中 $k=0,1,2,\ldots$。注意
$\mathfrak m_0^b(1)=\mathfrak m(e^b)$。
\end{definition}

于是容易得到下面的命题。

\begin{proposition}{3.6.10}\label{prop:3-6-10}
$(C,\mathfrak m^b)$ 也是滤过 $A_\infty$ 代数。并且
$\mathfrak m_0^b=0$ 当且仅当 $b$ 满足 $\widehat d(e^b)=0$；
在这种情况下，$\mathfrak m_1^b$ 正是定义 3.6.6 中的
$\delta_{b,b}$。
\end{proposition}

\begin{proof}
把 $\mathfrak m_k^b$ 延拓为 $\widehat B(C[1])$ 上的余导子
$\widehat d^{\,b}$，其公式为
\[
\begin{aligned}
&\widehat d^{\,b}(x_1,\ldots,x_k)\\
&\quad={}
\sum_{\ell=1}^{k+1}\sum_{\ell'=\ell}^{k+1}
(-1)^{\deg x_1+\cdots+\deg x_{\ell-1}+\ell-1}
x_1\otimes\cdots\otimes x_{\ell-1}\\
&\qquad{}\otimes
\mathfrak m(e^b x_\ell e^b\cdots e^b x_{\ell'-1}e^b)
\otimes x_{\ell'}\otimes\cdots\otimes x_k.
\end{aligned}
\]
当 $\ell'=\ell$ 时，上式右端相应位置的
$\mathfrak m(e^b x_\ell e^b\cdots e^b x_{\ell'-1}e^b)$
改用 $\mathfrak m(e^b)$。把这个公式与
\[
  \widehat d\circ\widehat d
  (e^b x_1e^b\cdots e^b x_ke^b)=0
\]
比较，便得到 $\widehat d^{\,b}\circ\widehat d^{\,b}=0$。
此外，注意恒等式
\[
  \widehat d(e^b)
  =e^b\mathfrak m(e^b)e^b
  =e^b\mathfrak m_0^b(1)e^b,
\]
便也证明了第二个断言。
\end{proof}

\begin{definition}{3.6.11}\label{def:3-6-11}
\liftIndex{I106}
对 $b_1,b_0\in\widehat{\mathcal M}(C)$，置
\[
  HF(C,b_1,b_0;\Lambda_{0,\mathrm{nov}})
  =H(C,\delta_{b_1,b_0}).
\]
当 $b_1=b_0$ 时，简写为
\[
  HF(C,b;\Lambda_{0,\mathrm{nov}})=H(C,\mathfrak m_1^b).
\]
\end{definition}

回忆，在 $A_\infty$ 代数 $C$ 中考虑同伦幺元时，我们曾像 (3.3.1)
那样把 $C$ 扩张为
\[
  C^+=C\oplus\Lambda_{0,\mathrm{nov}}\mathbf e^+
      \oplus\Lambda_{0,\mathrm{nov}}\mathbf f.
\]
下面研究 $C$ 与 $C^+$ 的定界余链之间的关系。

\begin{proposition}{3.6.12}\label{prop:3-6-12}
若 $b\in C[1]^0=C^1$ 是 $(C,\mathfrak m)$ 的定界余链，则它也是
$(C^+,\mathfrak m^+)$ 的定界余链。自然包含
\[
  (C,\mathfrak m^b)\longrightarrow(C^+,\mathfrak m^{+b})
\]
诱导同构
\[
  H(C,\mathfrak m_1^b)\longrightarrow
  H(C^+,\mathfrak m_1^{+b}).
\]
\end{proposition}

\begin{remark}{3.6.13}\label{rem:3-6-13}
由定理 4.2.45 将可看出，第二个断言蕴含自然包含
\[
  (C,\mathfrak m^b)\longrightarrow(C^+,\mathfrak m^{+b})
\]
给出这两个滤过 $A_\infty$ 代数之间的同伦等价。$A_\infty$ 代数之间
同伦等价概念的定义见第 4.2 节。
\end{remark}

\begin{proof}
第一个断言显然。下面证明第二个断言。自然包含显然是单射上链映射。
商上链复形由 $[\mathbf e^+]$ 和 $[\mathbf f]$ 生成；(3.3.5.2) 蕴含
\[
  d[\mathbf f]=[\mathbf e^+].
\]
故商复形无上同调，命题得证。
\end{proof}

现在把上述讨论应用于几何情形。对相对 spin 拉格朗日子流形 $L$，我们有
与之关联的滤过 $A_\infty$ 代数
$C(L;\Lambda_{0,\mathrm{nov}})$。

\begin{definition}{3.6.14}\label{def:3-6-14}
设
\[
  b\in C(L;\Lambda_{0,\mathrm{nov}})[1]^0
  =C^1(L;\Lambda_{0,\mathrm{nov}})
\]
上式两个空间之间的关系取等号
\translatornote{原书在这里写作第二个“$\in$”；但两边分别是
同一移位分次的两种写法，故改为等号。}
若 $\widehat d(e^b)=0$，就称 $b$ 是定界余链。
把所有定界余链 $b$ 构成的集合记作
$\widehat{\mathcal M}(L;J,\mathfrak s)$。若
$\widehat{\mathcal M}(L;J,\mathfrak s)$ 非空，就称 $L$
无障碍\liftIndex{I069}\liftIndex{I109}\liftIndex{I110}。
\end{definition}

解集 $\widehat{\mathcal M}(L;J,\mathfrak s)$ 可能依赖于给定的相容殆复结构
$J$ 以及多值截面系 $\mathfrak s$。不过，$L$ 的无障碍性与
$J,\mathfrak s$ 无关；这将在第 4.6 节证明。不会混淆时，为简便起见把
$\widehat{\mathcal M}(L;J,\mathfrak s)$ 写成
$\widehat{\mathcal M}(L)$。

\subsection{定界余链存在性的障碍}
\label{subsec:3-6-2}

上一小节引入了定界余链 $b$，并定义了上同调群
$H(C,\mathfrak m_1^b)$。在几何应用中，知道
$\widehat{\mathcal M}(L)$ 何时非空十分重要。本小节用
$(\overline C,\overline{\mathfrak m}_1)$ 的上同调，给出
$\widehat{\mathcal M}(C)$ 非空的条件；当
$C=C(L;\Lambda_{0,\mathrm{nov}})$ 时，这就是 $L$ 在
$\mathbb Q$ 上的通常上同调。为此，我们用代数方式在
$H(\overline C,\overline{\mathfrak m}_1)$ 中定义一列上同调类，称为障碍类。

令 $G\subset\mathbb R_{\geq0}\times2\mathbb Z$ 是满足条件 3.1.6 的
子幺半群。把它写成
\[
  G=\{(\lambda_i,\mu_i)\mid i=0,1,2,\ldots\},
  \qquad i<j\Longrightarrow\lambda_i\leq\lambda_j.
\]
（注意 $\lambda_0=\mu_0=0$ 且 $0<\lambda_1$。）考虑一个
$G$-有隙滤过 $A_\infty$ 代数 $C$，见定义 3.2.26。

\begin{remark}{3.6.15}\label{rem:3-6-15}
若 $G'$ 是另一个满足条件 3.1.6 且包含 $G$ 的子幺半群，则每个
$G$-有隙滤过 $A_\infty$ 代数也是 $G'$-有隙的。
\end{remark}

\begin{definition}{3.6.16}\label{def:3-6-16}
令 $G'\supseteq G$。若存在形如
\[
  b=\sum_{i=1}^{\infty}
  T^{\lambda_i'}e^{\mu_i'/2}b_i
\]
的定界余链，其中
\[
  b_i\in\overline C[1]^{-\mu_i'}
  =\overline C^{1-\mu_i'},
  \qquad
  G'=\{(\lambda_i',\mu_i')\mid i=1,2,\ldots\},
\]
就称一个 $G$-有隙滤过 $A_\infty$ 代数\emph{在 $G'$ 上无障碍}。
\end{definition}

\begin{lemma}{3.6.17}\label{lem:3-6-17}
$G$-有隙滤过 $A_\infty$ 代数无障碍，当且仅当它在某个
$G'\supseteq G$ 上无障碍。
\end{lemma}

这个引理的证明是显然的。现在给出本小节的主要结果。

\begin{theorem}{3.6.18}\label{thm:3-6-18}
像上面一样，设
\[
  G=\{(\lambda_i,\mu_i)\mid i=0,1,2,\ldots\},
  \qquad
  i<j\Longrightarrow\lambda_i\leq\lambda_j,
  \qquad \lambda_0=\mu_0=0,
\]
并令 $(C,\mathfrak m)$ 是一个 $G$-有隙滤过 $A_\infty$ 代数。
则存在闭上链序列 $o_k$ 和上链序列 $b_k$，具有下列性质。

\smallskip
\begin{equation}
  [o_k]\in H^{2-\mu_k}(\overline C,\overline{\mathfrak m}_1),
  \qquad b_k\in\overline C^{1-\mu_k}.
  \tag{3.6.19.1}\label{eq:3-6-19-1}
\end{equation}

\noindent\textup{(3.6.19.2)}
若所有满足 $\lambda_j<\lambda_k$ 的 $j$ 所对应的 $b_j,o_j$ 都已定义，
则闭上链 $o_k$ 有定义；上同调类 $[o_k]$ 依赖于这些 $b_j,o_j$。

\noindent\textup{(3.6.19.3)}
若 $[o_k]=0$，则上链 $b_k$ 有定义。

\noindent\textup{(3.6.19.4)}
$(C,\mathfrak m)$ 在 $G$ 上无障碍，当且仅当可以归纳地选取 $b_k$，
使所有 $[o_k]$ 都为零。
\end{theorem}

称 $[o_k]$（$k=1,2,\ldots$）为 $(C,\mathfrak m)$ 的\emph{障碍类}
\liftIndex{I107}\liftIndex{I108}\liftIndex{I349}。

\begin{proof}
由于 $(C,\mathfrak m)$ 是 $G$-有隙滤过 $A_\infty$ 代数，可以像定义
3.2.26 那样写成
\[
  \mathfrak m_k
  =\sum_{(\lambda_i,\mu_i)\in G}
   T^{\lambda_i}e^{\mu_i/2}\mathfrak m_{k,i}.
\]
这里
\[
  \mathfrak m_{k,i}\colon B_k\overline C[1]\longrightarrow\overline C[1],
  \qquad \mathfrak m_{k,0}=\overline{\mathfrak m}_k.
\]
为了处理若干 $\lambda_i$ 相等的情形，先准备一些记号。若
$\lambda_i$ 是集合 $\{\lambda_j\}$ 中第 $k$ 小的非零数，就写
$\lambda_i=\lambda_i^{(k)}$，并置
$\lambda^{(k)}=\lambda_i^{(k)}$。例如，若
\[
  0<\lambda_1=\cdots=\lambda_\ell<\lambda_{\ell+1},
\]
就写
\[
  \lambda_1^{(1)}=\lambda_1,\ldots,
  \lambda_\ell^{(1)}=\lambda_\ell,
  \qquad
  \lambda_{\ell+1}^{(2)}=\lambda_{\ell+1}.
\]
下面关于 $(k)$ 归纳地构造 $o_i$ 和 $b_i$。首先考虑满足
\[
  \lambda_1^{(1)}=\cdots=\lambda_\ell^{(1)}
  <\lambda_{\ell+1}^{(2)},
\]
的情形，并置
\[
  o_i=\mathfrak m_{0,i}(1)
  \in\overline C[1]^{1-\mu_i}
  =\overline C^{2-\mu_i},
  \qquad i=1,\ldots,\ell.
\]
注意 $\mathfrak m_{0,0}=0$，于是
\[
\begin{aligned}
  T^{\lambda^{(1)}}
  \sum_{i=1}^{\ell}e^{\mu_i/2}
  \mathfrak m_{1,0}(o_i)
  &=\sum_{i=1}^{\ell}T^{\lambda_i}e^{\mu_i/2}
    \mathfrak m_{1,0}\mathfrak m_{0,i}(1)\\
  &\equiv\mathfrak m_1(\mathfrak m_0(1))
  \pmod{T^{\lambda^{(1)}}\Lambda^+_{0,\mathrm{nov}}}.
\end{aligned}
\]
由 $A_\infty$ 公式，末项为零。又由于
$1\leq i\ne i'\leq\ell$ 时 $\mu_i\ne\mu_{i'}$，并且
$\mathfrak m_{1,0}=\overline{\mathfrak m}_1$，所以对
$i=1,\ldots,\ell$ 都有
$\overline{\mathfrak m}_1(o_i)=0$。因此，$o_i$ 定义了上同调类
\[
  [o_i]\in H^{2-\mu_i}
  (\overline C,\overline{\mathfrak m}_1),
  \qquad i=1,\ldots,\ell.
\]
若这些 $[o_i]$ 全都为零，就可以选取
$b_i\in\overline C^{1-\mu_i}$，使
\[
  \mathfrak m_{1,0}(b_i)+o_i=0.
\]
进一步，定义 $b_{(1)}$ 如下：
\[
  b_{(1)}=\sum_{i=1}^{\ell}
  T^{\lambda_i}e^{\mu_i/2}b_i\in C^1,
\]
由此得到
\begin{equation}
\begin{aligned}
  \mathfrak m(e^{b_{(1)}})
  &\equiv\sum_{i=1}^{\ell}
  \bigl(
    T^{\lambda_i}e^{\mu_i/2}\mathfrak m_{0,i}(1)
    +T^{\lambda_i}e^{\mu_i/2}\mathfrak m_{1,0}(b_i)
  \bigr)\\
  &=0
  \pmod{T^{\lambda^{(1)}}\Lambda^+_{0,\mathrm{nov}}}.
\end{aligned}
\tag{3.6.20}\label{eq:3-6-20}
\end{equation}

现在归纳地假设：对满足 $\lambda_i<\lambda^{(k)}$ 的 $i$，已经选取
$b_i$，使得
\begin{equation}
  \mathfrak m(e^{b_{(k-1)}})
  \equiv0
  \pmod{T^{\lambda^{(k-1)}}\Lambda^+_{0,\mathrm{nov}}},
  \tag{3.6.21}\label{eq:3-6-21}
\end{equation}
其中
\begin{equation}
  b_{(i)}=\sum_{j=1}^{*_{i}}
  T^{\lambda_j}e^{\mu_j/2}b_j\in C^1,
  \qquad
  *_{i}=\max\{j\mid\lambda_j<\lambda^{(i+1)}\}.
  \tag{3.6.22}\label{eq:3-6-22}
\end{equation}

\begin{definition}{3.6.23}\label{def:3-6-23}
当以下关系成立时：
\[
  \lambda_n^{(k)}=\cdots=\lambda_{n+m}^{(k)}
  <\lambda_{n+m+1}^{(k+1)}
\]
此时 $*_{k-1}=n-1$；对每个 $h=n,\ldots,n+m$，把
$\mathfrak m(e^{b_{(k-1)}})$ 中
$T^{\lambda_h}e^{\mu_h/2}$ 的系数记作
\[
  o_h\in\overline C^{2-\mu_h}.
\]
\end{definition}

\begin{lemma}{3.6.24}\label{lem:3-6-24}
$\overline{\mathfrak m}_1(o_h)=0$。
\end{lemma}

\begin{proof}
再次应用 $A_\infty$ 公式，得到
\[
\begin{aligned}
0
&=(\mathfrak m\circ\widehat d)(e^{b_{(k-1)}})\\
&\equiv
T^{\lambda^{(k)}}\mathfrak m
\left(\sum_{h=n}^{n+m}e^{\mu_h/2}o_h\right)\\
&\equiv
\sum_{h=n}^{n+m}T^{\lambda_h}e^{\mu_h/2}
\mathfrak m_{1,0}(o_h)
\pmod{T^{\lambda^{(k)}}\Lambda^+_{0,\mathrm{nov}}}.
\end{aligned}
\]
由此，对每个 $h=n,\ldots,n+m$ 都有
$\mathfrak m_{1,0}(o_h)=0$，引理得证。
\end{proof}

至此已经定义了上同调类 $[o_h]$。若它们全都为零，可以选取
$b_h\in\overline C^{1-\mu_h}$，使其满足
\begin{equation}
  \mathfrak m_{1,0}(b_h)+o_h=0.
  \tag{3.6.25}\label{eq:3-6-25}
\end{equation}

\begin{lemma}{3.6.26}\label{lem:3-6-26}
\[
  \mathfrak m(e^{b_{(k)}})
  \equiv0
  \pmod{T^{\lambda^{(k)}}\Lambda^+_{0,\mathrm{nov}}},
\]
其中 $b_{(k)}$ 由~\eqref{eq:3-6-22} 定义。
\end{lemma}

\begin{proof}
由于 $G=\{(\lambda_i,\mu_i)\mid i=0,1,2,\ldots\}$ 是幺半群，
\[
  \mathfrak m(e^{b_{(k-1)}})
  \equiv
  \sum_{h=n}^{n+m}T^{\lambda_h}e^{\mu_h/2}o_h
  \pmod{T^{\lambda^{(k)}}\Lambda^+_{0,\mathrm{nov}}}.
\]
引理遂由~\eqref{eq:3-6-25} 得出。
\end{proof}

至此已经证明 (3.6.19.1)、(3.6.19.2) 和 (3.6.19.3)。容易看出，
若所有 $b_k$ 都已定义（因而所有 $[o_k]$ 都为零），则
\begin{equation}
  b=\sum_kT^{\lambda_k}e^{\mu_k/2}b_k
  =\lim_{k\to\infty}b_{(k)}
  \tag{3.6.27}\label{eq:3-6-27}
\end{equation}
是定界余链。事实上，由引理 3.6.26，对任意 $k$ 都有
\[
  \mathfrak m(e^b)\equiv0
  \pmod{T^{\lambda_k}\Lambda^+_{0,\mathrm{nov}}},
\]
故 $\mathfrak m(e^b)=0$。因而
\[
  \widehat d(e^b)=e^b\mathfrak m(e^b)e^b=0.
\]
反过来，若 $(C,\mathfrak m)$ 在 $G$ 上无障碍，就能找到形如
\eqref{eq:3-6-27} 的定界余链 $b$。于是 $b_k$ 满足
\eqref{eq:3-6-21} 和~\eqref{eq:3-6-25}，也就是说，可以定义
$b_k$ 使所有 $[o_k]$ 都消失。定理 3.6.18 证毕。
\end{proof}

\begin{remark}{3.6.28}\label{rem:3-6-28}
$(C,\mathfrak m)$ 可能在某个 $G$ 上有障碍，却在某个更大的 $G'$ 上无障碍。
这一现象同模空间 $\widehat{\mathcal M}(C)$ 的有理点的存在性有关；
可比较\textbf{[Fuk03I]}。
\end{remark}

若像定理 3.5.11 那样
$(C,\mathfrak m)=(C(L;\Lambda_{0,\mathrm{nov}}),\mathfrak m)$，就把
$o_k$ 写成 $o_k(L)$。把 $G=G(L)$ 的元素枚举为
\[
\begin{split}
\bigl\{(\lambda(\beta_i),\mu(\beta_i))\mathrel{\big|}
  &\ \beta_0=0,\quad
  \beta_1,\beta_2,\ldots\\
  &\in\{\beta\in\Pi(L)\mid\mathcal M(\beta)\ne\varnothing\}\bigr\},
\end{split}
\]
并使 $\lambda(\beta_i)\leq\lambda(\beta_{i+1})$。上式的内层集合按唯一
可成立的集合构造排印
\translatornote{原书在内层集合的 $\Pi(L)$ 后多排了一个右括号；译文删去。}。
\mbox{其中}
\[
  \lambda(\beta_i)=\omega[\beta_i],
  \qquad
  \mu(\beta_i)=\mu_L(\beta_i)
\]
分别是能量（辛面积）和马斯洛夫指标；并置
\[
  o(L;\beta_i):=o_i(L).
\]
